\documentclass[10pt,a4paper]{amsart}
\usepackage{amsmath,amssymb,amsthm,graphicx}
% Retained minimal preamble for the Stacks Project chapter more-morphisms.tex (More on Morphisms).
% Verbatim excerpt of the chapter preamble: Xy-pic setup (native preamble lines 14-18) and the native
% theorem declarations with the mathematical macros (native preamble lines 154-217). Copied mechanically;
% nothing added, removed or redefined. The wrapper class and the amsmath/amssymb/amsthm/graphicx base
% packages are supplied by the assembler, so the class/externaldocument part of the native preamble is omitted.
% For commutative diagrams we use Xy-pic
\usepackage[all]{xy}

% We use 2cell for 2-commutative diagrams.
\xyoption{2cell}
% Theorem environments.
%
\theoremstyle{plain}
\newtheorem{theorem}[subsection]{Theorem}
\newtheorem{proposition}[subsection]{Proposition}
\newtheorem{lemma}[subsection]{Lemma}

\theoremstyle{definition}
\newtheorem{definition}[subsection]{Definition}
\newtheorem{example}[subsection]{Example}
\newtheorem{exercise}[subsection]{Exercise}
\newtheorem{situation}[subsection]{Situation}

\theoremstyle{remark}
\newtheorem{remark}[subsection]{Remark}
\newtheorem{remarks}[subsection]{Remarks}

\numberwithin{equation}{subsection}

% Macros
%
\def\lim{\mathop{\mathrm{lim}}\nolimits}
\def\colim{\mathop{\mathrm{colim}}\nolimits}
\def\Spec{\mathop{\mathrm{Spec}}}
\def\Hom{\mathop{\mathrm{Hom}}\nolimits}
\def\Ext{\mathop{\mathrm{Ext}}\nolimits}
\def\SheafHom{\mathop{\mathcal{H}\!\mathit{om}}\nolimits}
\def\SheafExt{\mathop{\mathcal{E}\!\mathit{xt}}\nolimits}
\def\Sch{\mathit{Sch}}
\def\Mor{\mathop{\mathrm{Mor}}\nolimits}
\def\Ob{\mathop{\mathrm{Ob}}\nolimits}
\def\Sh{\mathop{\mathit{Sh}}\nolimits}
\def\NL{\mathop{N\!L}\nolimits}
\def\CH{\mathop{\mathrm{CH}}\nolimits}
\def\proetale{{pro\text{-}\acute{e}tale}}
\def\etale{{\acute{e}tale}}
\def\QCoh{\mathit{QCoh}}
\def\Ker{\mathop{\mathrm{Ker}}}
\def\Im{\mathop{\mathrm{Im}}}
\def\Coker{\mathop{\mathrm{Coker}}}
\def\Coim{\mathop{\mathrm{Coim}}}

% Boxtimes
%
\DeclareMathSymbol{\boxtimes}{\mathbin}{AMSa}{"02}

%
% Macros for moduli stacks/spaces
%
\def\QCohstack{\mathcal{QC}\!\mathit{oh}}
\def\Cohstack{\mathcal{C}\!\mathit{oh}}
\def\Spacesstack{\mathcal{S}\!\mathit{paces}}
\def\Quotfunctor{\mathrm{Quot}}
\def\Hilbfunctor{\mathrm{Hilb}}
\def\Curvesstack{\mathcal{C}\!\mathit{urves}}
\def\Polarizedstack{\mathcal{P}\!\mathit{olarized}}
\def\Complexesstack{\mathcal{C}\!\mathit{omplexes}}
% \Pic is the operator that assigns to X its picard group, usage \Pic(X)
% \Picardstack_{X/B} denotes the Picard stack of X over B
% \Picardfunctor_{X/B} denotes the Picard functor of X over B
\def\Pic{\mathop{\mathrm{Pic}}\nolimits}
\def\Picardstack{\mathcal{P}\!\mathit{ic}}
\def\Picardfunctor{\mathrm{Pic}}
\def\Deformationcategory{\mathcal{D}\!\mathit{ef}}
\usepackage{cleveref}
\crefname{originallemma}{original lemma}{original lemmas}
\Crefname{originallemma}{Original lemma}{Original lemmas}
\makeatletter
\newlabel{lemma-thickening-affine-scheme}{{06AD}{source}{}{sourcepointer-1}{}}
\newlabel{lemma-thickening-affine-scheme@cref}{{[originallemma][0][]06AD}{[1][0][]source}{}{}{}}
\makeatother
\crefname{originallemma}{original lemma}{original lemmas}
\Crefname{originallemma}{Original lemma}{Original lemmas}
\makeatletter
\newlabel{limits-lemma-ample-on-reduction}{{09MW}{source}{}{sourcepointer-2}{}}
\newlabel{limits-lemma-ample-on-reduction@cref}{{[originallemma][0][]09MW}{[1][0][]source}{}{}{}}
\makeatother
\crefname{originallemma}{original lemma}{original lemmas}
\Crefname{originallemma}{Original lemma}{Original lemmas}
\makeatletter
\newlabel{morphisms-definition-relatively-ample}{{01VH}{source}{}{sourcepointer-3}{}}
\newlabel{morphisms-definition-relatively-ample@cref}{{[originallemma][0][]01VH}{[1][0][]source}{}{}{}}
\makeatother
\title{Relative ampleness of an invertible sheaf is insensitive to thickening the source and the base}
\author{The Stacks Project Authors}
\date{stacks-project:more-morphisms.tex:local-snapshot-20261011; GNU FDL 1.2 or later}
\begin{document}
\maketitle
\noindent\textit{Editorial note.} All delivered mathematical text is the author's own, copied from the pinned Stacks Project chapter file more-morphisms.tex; no statement or proof step was rewritten, reordered or removed.\par\medskip
\noindent\textit{Editorial note.} Preamble provenance, as required for this source: the first line of the pinned chapter file is \textbackslash{}input\{preamble\}, and the native preamble was not staged with the source, so no preamble text lives inside the chapter. The retained preamble used here is the project's single Stacks-wide preamble (the same 217-line / 8116-byte file that the already published chow and etale-cohomology cards use), copied mechanically line for line from sources/stacks-chow/preamble.tex lines 14-18 (Xy-pic setup) and 154-217 (theorem declarations and mathematical macros). Not one character was changed, no package was added, no macro was defined, and the excerpt serves only as preamble, never as a span coordinate.\par\medskip
\noindent\textit{Editorial note.} The same-chapter pointers are located at their \textbackslash{}label lines inside the pinned chapter file; the cross-chapter pointers are located in the project's own tag list staged in the snapshot (sources/stacks-chow/tags.txt), at the verbatim lines carrying their canonical tags.\par\medskip
\noindent\textit{Allowed standard prerequisites (editorial).}\begin{itemize}
\item Relative ampleness is checked on affine opens of the base. Morphisms, Definition labelled morphisms-definition-relatively-ample; invoked at lines 321-323. Reduces the relative statement to affine bases.
\item Affine opens of S correspond to affine opens of the thickening S'. Same chapter, Lemma labelled lemma-thickening-affine-scheme; invoked at lines 324-326. Transports the affine-open reduction across the thickening and identifies the restricted situation with (X subset X') over (S subset S').
\item Ampleness is insensitive to passing from a thickening to its reduction. Limits, Lemma labelled limits-lemma-ample-on-reduction; invoked at line 331. Supplies the affine case of the equivalence.
\end{itemize}\medskip
\ifdefined\hypertarget\hypertarget{sourcepointer-1}{}\fi
\noindent\textit{Source reference.} Original lemma 06AD: sources/stacks-more-morphisms/pinned-original-source.tex, source lines 157--157. Navigation pointer for \textbackslash{}ref\{lemma-thickening-affine-scheme\} used inside the counted span (line 324); the argument that consumes it is retained in the span.\par\medskip
\ifdefined\hypertarget\hypertarget{sourcepointer-2}{}\fi
\noindent\textit{Source reference.} Original lemma 09MW: sources/stacks-chow/tags.txt, source lines 12194--12194. Navigation pointer for \textbackslash{}ref\{limits-lemma-ample-on-reduction\} used inside the counted span (line 331); the argument that consumes it is retained in the span.\par\medskip
\ifdefined\hypertarget\hypertarget{sourcepointer-3}{}\fi
\noindent\textit{Source reference.} Original lemma 01VH: sources/stacks-chow/tags.txt, source lines 2446--2446. Navigation pointer for \textbackslash{}ref\{morphisms-definition-relatively-ample\} used inside the counted span (line 323); the argument that consumes it is retained in the span.\par\medskip
\begin{lemma}
\label{lemma-thicken-property-relatively-ample}
Let $(f, f') : (X \subset X') \to (S \subset S')$ be a morphism
of thickenings. Let $\mathcal{L}'$ be an invertible sheaf on $X'$
and denote $\mathcal{L}$ the restriction to $X$.
Then $\mathcal{L}'$ is $f'$-ample if and only if
$\mathcal{L}$ is $f$-ample.
\end{lemma}

\begin{proof}
Recall that being relatively ample is a condition for each
affine open in the base, see
Morphisms, Definition \ref{morphisms-definition-relatively-ample}.
By Lemma \ref{lemma-thickening-affine-scheme}
there is a 1-to-1 correspondence between
affine opens of $S$ and $S'$.
Thus we may assume $S$ and $S'$ are affine
and we reduce to proving that
$\mathcal{L}'$ is ample if and only if
$\mathcal{L}$ is ample.
This is Limits, Lemma \ref{limits-lemma-ample-on-reduction}.
\end{proof}

\end{document}
